\chapter{Source Acquisition and Retrieval Pipeline}

This chapter describes how official UniTrento sources are collected, stored, transformed into searchable units, and indexed for retrieval. These components provide the controlled document base used in the experimental comparison between standard BM25 retrieval and BM25 retrieval supported by LLM-based query rewriting.

\section{Source Repository}
The repository contains official institutional material related to student bureaucracy, including programme pages, university regulations, yearly calls, administrative FAQs, and PDF guidelines. Sources are restricted to trusted UniTrento, DISI, CLA, and Opera Universitaria domains.

The final experiment uses two overlapping corpus views over the same stored repository:

\begin{itemize}
    \item \textbf{Polished corpus}: a manually curated set of 27 documents selected to cover the benchmark topics with limited unrelated material;
    \item \textbf{Noisy corpus}: a broader set of 463 current documents obtained through wider crawling and discovery.
\end{itemize}

Every polished document is also part of the noisy corpus. The two views are represented through exact source-set membership rather than through separate repositories or separate indexes. This keeps document identity, extraction, and indexing fixed while allowing retrieval to be filtered to one corpus or the other.

\section{Storage and Versioning}
Each source is identified by a stable document identifier derived from its canonical URL. Revisions are stored separately so that changed content can be tracked without losing provenance. The current revision of each document is the one materialized for extraction and indexing.

This separation is important for yearly regulations and calls, where multiple versions may coexist. It also allows the benchmark corpus membership to be changed without crawling or extracting the same document again.

\section{Extraction}
Raw HTML and PDF sources are transformed into smaller searchable units. HTML pages are extracted into anchored sections or subsections. PDF documents are extracted at page level. Each segment retains provenance metadata, including the document and revision identifiers, source URL, title, anchor type, anchor value, and segment order.

Segment-level extraction makes the returned evidence easier to inspect and judge than whole-document retrieval. It also allows links to point to the relevant HTML heading or PDF page.

\section{BM25 Indexing}
Each extracted segment is materialized as one BM25 record and indexed in OpenSearch. Title and heading fields receive stronger weights than body text because matches in structural fields are usually more informative than isolated matches in long passages.

The index contains the union of all current source documents. Exact corpus selection is applied at query time through the indexed source-set identifiers. Therefore, switching between the polished and noisy corpora does not require another crawl, extraction pass, or index rebuild.

For the thesis retrieval profile, results are collapsed to one representative hit per document. This prevents several segments from the same file from occupying the complete top-$k$ window and makes the evaluation closer to source selection rather than passage duplication.

\section{Evaluation Dataset}
The final benchmark contains 26 Italian student-style questions covering topics such as study plans, internships, graduation, language requirements, scholarships, enrolment, taxes, and international mobility.

For each query, the system retrieves the top three distinct source documents. The same query set is used in every condition. Expected sources are retained as discovery aids, but final evaluation is based on manual relevance judgements of the retrieved results rather than on automatic URL matching.

\section{Experimental Retrieval Pipeline}
The experiment compares original and rewritten queries while keeping extraction, BM25 configuration, ranking logic, and result depth fixed.

\begin{figure}[h]
\centering
\fbox{\begin{minipage}{0.90\linewidth}
\centering
Stored sources $\rightarrow$ current revisions $\rightarrow$ HTML/PDF extraction $\rightarrow$ BM25 index\\[0.6em]
Original query $\rightarrow$ BM25 $\rightarrow$ top-3 distinct documents\\
Original query $\rightarrow$ two-layer LLM rewrite $\rightarrow$ BM25 $\rightarrow$ top-3 distinct documents
\end{minipage}}
\caption{Simplified architecture and experimental comparison.}
\label{fig:pipeline}
\end{figure}

Two corpus filters are evaluated: polished and noisy. Three independently generated rewrite sets are frozen and then reused unchanged on both corpora. Reusing the same rewrites is necessary to separate rewrite variation from the effect of corpus noise.

The thesis measures retrieval only. In the broader product prototype, retrieved segments can later be passed to a generative model to produce a grounded answer with citations, but answer generation is outside the core experiment.
\par
